% Created 2024-03-15 ven. 23:40
% Intended LaTeX compiler: pdflatex
\documentclass[11pt]{article}
\usepackage[utf8]{inputenc}
\usepackage[T1]{fontenc}
\usepackage{graphicx}
\usepackage{longtable}
\usepackage{wrapfig}
\usepackage{rotating}
\usepackage[normalem]{ulem}
\usepackage{amsmath}
\usepackage{amssymb}
\usepackage{capt-of}
\usepackage{hyperref}
\date{}
\title{discours présentation R-GTMG}

\begin{document}

\maketitle

Hello, today I present my works about reevaluation in rule-based graph
transformation modeling systems.

\section{Introduction}
\label{sec:orgded3f04}

Modeling complex 3D objects is tedious process often with many operations.

Nowadays, all modeling kernels record these process' in what's called a
parametric specification to modify and reevaluate them.

One can modify topological and geometric parameters and even add, delete or move
operations.

Let's see a specification it will be our guiding thread throughout this presentation:
\begin{itemize}
	\item First operation creates the square face f1 and its edges including e1.
	\item Next one extrudes f1 into a cube and e1 into the face f2.
	\item Triangulating f2 creates several faces including f7.
	\item Last operation colors f7.
\end{itemize}

Now let's modify this specification: we add the insertion of a vertex on edge e1
and we reevaluate.

\begin{itemize}
	\item f1 and e1 are identical
	\item the vertex splits e1 into e5 and e6
	\item The extrusion applies on f1, this time e5 and e6 are extruded into f3
	      and f4.
\end{itemize}

Finally, how can we continue reevaluating since f2 is no longer here ? The same goes
for f7.

Intuitively, f2 is split into f3 and f4, we want to triangulate f3 and/or f4.

Right now, that's when we need a persistent naming scheme to identify cells'
histories to be able to say, for example, ``f3 and f4 are built like f2,
then we triangulate these two faces.''

% Maintenant que la problématique est posée, on va voir le plan.
Now the issue is raised, let's see the outline.

\section*{Outline}
\label{sec:orgdf64a4e}
\begin{itemize}
	\item We aim to allow the reevaluation of modified specifications especially
	      by adding, deleting and moving operations.
	\item To this end we will see two formalisms for objets and operations
	      % \begin{itemize}
	      %   \item Celui des cartes généralisées pour les objets
	      %   \item Celui des règles de transformation de graphe Jerboa pour les
	      %         opérations de modélisation
	      % \end{itemize}
	\item Our contributions for a persistent naming scheme:
	      % \begin{itemize}
	      %   \item d'une part un mécanisme de nommage persistant basé sur les
	      %         règles
	      %   \item d'autre part un mécanisme de réévaluation exploitant ce
	      %         nommage persistant
	      % \end{itemize}
	\item And then I will conclude.
\end{itemize}

\section{G-maps}
\label{sec:orgf33c57b}

We start with generalized maps, or G-maps, to represent objects

Consider this object

It can be decomposed into two volumes connected with 3-links in green

Its volumes are faces connected by 2-links in blue

Which are edges connected by 1-links in red

And finally, these edges are vertices connected by 0-links in black.

This G-map is a graph whose nodes and arcs are called darts and links.

\section{Orbits}
\label{sec:org99b57a4}

G-maps have the ability to represent cells with sub-graphs called orbits

For example~:

The vertex incident to dart ``a'' contains ``a'', all the darts reachable
through 1,2,3-links and the links themselves.

This edge contains ``a'' and darts reachable through 0,2,3-links.

A face contains darts reachable through 0,1,3-links.

A volume contains those reachable through 0,1,2-links.

Orbits can also represent all kind of topological, not limited to cells, such as
this volume edge.

\section{Jerboa graph transformation rules}
\label{sec:org89fe9a5}

We have formalized the objects with G-maps not let's get to the operations with
Jerboa graph transformation rules.

Here is the triangulation of a face.

Generally speaking, a rule has a left membre used to match a sub-graph in the
object to transform and a right member which describe the transformation.

On the left, the node n0 holds the face <0,1,3> orbit type.

When applying the rule on dart ``a0'', it matches all the face incident to
``a0'' (in yellow).

On the right, the node n0 is preserved thus the matched darts are preserved. %
The rewriting of this implicit 1-arc into an underscore in n0 denotes the
deletion of the 1-links in the object. %

n2 is a created node, as such, the darts it creates are copied from those
matched by n0. %
In the rule, 0-links are rewritten into 1-links and the 1-links into 2-links. %
In the object, the vertex, in blue, is created as the dual of the face, in
yellow, by replacing 0-,1-links with 1-,2-links. %

Finally, n1 is created and presents both deletion and rewriting of links. %

In the rule, an explicit 1-arc connects n0 to n1 which links their darts 1-to-1 in the object. %
The same goes for this 0-link between n1 and n2. %

\subsection{Orbits}
\label{sec:orgcf72bd7}

We can also track the orbits.

Let's track the volume face incident to n0.

On the left, it contains n0 with its 0-,1-arcs.

On the right, it contains n0 with its 0-arc, the 1-arc connecting it to n1, the
0-arc connecting n1 to n2 and n2's 1-arc.

Note that there is only one face orbit in the rule yet there are 4 of them in
the object. %
That's because the second implicit link on the left is either deleted or
rewritten into 2 (out of the 0,1 orbit type). %
Such rewriting splits the face. %

\subsection{Origins}
\label{sec:org341560c}

Moreover, the implicit 0,1-arcs on the right are rewritten from the implicit 0-arc on the left. %
Each one of the faces on the right originates from a single edge on the left.

For example, this green face comes from this green edge.

A static analysis of the rules, performed at pre-compute, time makes it possible
to detect the events on orbits independently of the object onto which the
operations are applied.
We refer to these events as the split, merging and so on.

\section*{Outline: persistent naming}
\label{sec:orgaa5036e}

That's it the formalisms, let's get to the persistent with our parametric
specification.

\section{Persistent naming}
\label{sec:org84648da}

\subsection{Persistent naming: darts}
\label{sec:org53f1b92}

The history of each dart, namely the rules and nodes that create them, is
recorded and used as their persistent names. %

This rule, creating a square face, has 8 node including n3 and n6.

With this first application, n3 creates the dart 3, its history [1n3].

n6 creates the dart 6, its history is [1n6]

To extrude the square, we select the dart 6. However, its name may change. Thus,
when applying the rule, we make a snapshot of the dart's history and use it as a
persistent name for the specification.

n1, on the left, matches dart 3 while n2 and n4, on the right, create its copies 33 and 35. %
As such, we build their histories from dart 3.

Dart 33 history is [1n3;2n2]. %
For dart 35 it's [1n3;2n4]. %
Dart 3 history is update and becomes [1n3;2n1], the same goes for dart 6. %

Next we triangule the front face which we select through dart 33.
Again we snapshot the history into a persistent name.

Darts 33 and 35 are updated with n0 on the left.
For 35 this third application gives [1n3;2n4;3n0], the same goes for 33.

Finally, the face is colored.

Here, we have just given unique persistent names to darts based on their histories.

Yet, topological parameters are orbits and not darts.
Now we must represent the orbits histories.

\subsection{Persistent naming: orbits}
\label{sec:orgf5219a2}

Now, we use the history of a dart to reconstitute the history of an
orbit.

Let's consider the orbit designated with [1n3;2n4;3n0] in the color operation.
n0, in yellow, denotes a volume face, then we start tracing the history from a
face which is drawn here.

From the persistent name, we start from the last step being the third operation. %

As I told earlier, this face is split from a matched face and originates from an
edge. %
We transfer these informations into the history.

So, this face incident to n0, is tracked, the black arrow, back to a matched
face split by the triangulation. It originates, the red arrow, to from an edge.

We keep backtracking the history up to the second operation, the
extrusion and its node n4.

Now we track this face and this edge both incident to n4 respectively in blue
and yellow boxes.

The face is created through the extrusion of an edge and the edge itself is
created from a dart.

Finally, we backtrack to the first operation, the square creation which creates
these edge and dart.

We have reconstituted an orbit history solely from the persistent name of a dart
and the rules from the specification.

We did not need to reapply any of the rules.
% On fait ça uniquement à partir de la spécification et des règles, sans utiliser
% l'objet et sans réappliquer les règles.

This process is performed for each persistent name and only those.

In the end we only reconstitute orbit histories if there is a corresponding
persistent name in the specification.

\section*{Outline: reevaluation}
\label{sec:org04c39be}

Thanks to these persistent names, we can now proceed to the reevaluation.

\section{Reevaluation}
\label{sec:org3b38850}

Let's reevaluate the specification from the standpoint of the history we just
built.  %

The fist operation recreates the square identically. %
n3 has created a unique dart 3.

The next operation has been added, I miss the place and time to show it, still
its analysis reveals it split the matched edge. %
In the history we now trace a branch for each edge.

The extrusion creates a face from each edge, an edge from each dart. %
n4 creates several darts yet only one copy of 3, being 37, and one copy of 4,
being 52.

Finally we reapply the triangulation on each branch. %
Darts 37 and 52 remain identical. %
It was the last operation, darts 37 and 52 now designate the faces to color. %

This way, we can find all the topological entities eligible during the
reevaluation. %
The user has the choice to reapply operations for each branch. %

\section{Conclusion}
\label{sec:orga7a7cfd}

To conclude, we propose a reevaluation mechanism built for rule-based
modeling systems.

To this end, we designed a two layers persistent naming scheme:

One for the darts, making the naming unique;

The other of for the orbits, representing a complete history tracing the
evolution of an orbit up to its origins.

The naming of an orbit is solely built from the data obtain during the analysis of
the rules used in a specification. The analysis itself is statically performed
at pre-compute time.

As such, this approach distinguishes itself from other where all cells of a given
dimension (usually edges and faces) are being track throughout the design process.

During the reevaluation, we track all the possible histories allowing setting up strategies such as reapplying an operation on one or many operations.

We currently work on the integration of scripted rules involving usual control
structures such as loops and conditionals within this reevaluation mechanism.
\end{document}
